\scanpage{11}
\mspage{10}

On definit dans l'ens $\mathcal{F}$ des \uline{figures} \inserted{$F$}, comme l'ens des parties $\Phi={\Sigma}'_{F}$ de $\mathcal{F}_{el}$, qui sont a) des configurations

b) \unsure{fermées|fermees} pour la relation d'ordre

c) loc finies. Figures \uline{finies} : parties finies de $\mathcal{F}_{el}$ satisfaisant a) et b).

% slanted
\margin{Figures compatibles : qui le sont comme configurations. La \uline{réunion} \inserted{des \unsure{$E_{el}$|Ens}} supports sera alors la figure somme, et l'intersection sera aussi une figure -- la figure intersection \ldots}

Notons comme propriétés, en plus des propriétés a) : d) de $4^\circ$) ci-dessus, la relation entre \textcircled{2} et \textcircled{3}

\deleted{Axiome Si}

\uline{Fig}$_1$ Si $\deleted{\unsure{\Phi|F}}\unsure{X|F|\mathbf{X}}_1,\deleted{F}X_2\leq\deleted{F}X$ sont des fig. él. alors $\{\deleted{F}X_1,\deleted{F}X_2\}\in R$ i.e. $\{\deleted{F}X_1,\deleted{F}X_2\}$ est une configuration.

[Donc l'ens des $\deleted{F}X'\leq\deleted{F}X$ est une configuration, et en vertu de $4^\circ$ a) b) c'est $\overline{\mathrm{m}}$ une \uline{figure} : la figure \unsure{associée|associé} : $F$\unsure{\ldots|\,--}\ ]

% slanted
\margin{On l'appelle la figure élémentaire \unsure{associé|associée} à $F$}
